% ==============================================================================
% 第 13 章：第二部分小结 (Chapter 13: Summary)
% ==============================================================================

\chapter{第二部分小结}
\label{sec:ch13_summary}

本书的第二部分专注于\term{处理任意大小数据}（Arbitrarily Large Data）的程序设计。正如你所体会到的，当软件能够处理没有预设大小限制的海量信息时，其价值才真正得以释放。掌握“任意大小的数据”是迈向真正职业程序员道路上的决定性一步。在此，我们建议你牢记以下三条核心经验：

\begin{notebox}[第二部分核心收获]
\begin{enumerate}
  \item \textbf{自引用与交叉引用的设计配方细化}：
  本部分对设计配方进行了重大扩展，使其能够系统化地应对数据定义中的\textbf{自引用}与\textbf{交叉引用}。
  \begin{itemize}
    \item 数据定义中出现自引用，必然指导我们在函数模板与函数体中构建对应的\term{结构型递归调用}；
    \item 数据定义中出现对其他复杂数据定义的交叉引用，必然指导我们设计专门处理对应数据的\term{独立辅助函数}。
  \end{itemize}

  \item \textbf{分解与组合的设计范式 (Design by Composition)}：
  复杂问题必然要求分解为相互协作的子任务。当你分解一个问题时，需要两个支撑要素：分别解决各子任务的函数，以及用于组合这些独立解的数据定义。
  为了确保在耗费精力完成各个独立函数后整体系统能够严丝合缝地协同工作，你必须\textbf{严格制定包含明确函数签名、目的说明与接口契约的“愿望清单”（Wish List）}。
  
  在以下三种情形中，“分解—组合”的设计手法尤为关键：
  \begin{itemize}
    \item 问题描述中明确或隐含了不同专业维度的辅助任务；
    \item 函数体的编写需要遍历另一个任意大小的数据集合；
    \item 求解一个\textbf{更一般化、更通用的问题}往往比直接求解原始特定问题更加容易且优美。
  \end{itemize}

  \item \textbf{工程务实与领域知识 (Pragmatics Matter)}：
  设计交互式世界程序时，必须深入理解 \code{big-bang} 的各类事件子句及其语义边界；解决特定数学或工程问题时，必须清楚所选语言及其标准库提供了哪些现成的高性能算法与原语。
\end{enumerate}
\end{notebox}

虽然本部分主要以列表（List）作为任意大小数据的典型范例——因为它们在 Haskell、Lisp、ML、Racket 和 Scheme 等现代函数式语言中扮演着不可替代的角色——但这些设计精髓完全适用于所有类似的数据形态：文件系统、多级目录树、数据库表、关系网络等。

在本书的第四部分中，我们将进一步探索交织与层次更加复杂的“大型”结构化数据（如多叉树、S-表达式与 XML 文档），并展示设计配方如何无缝缩放到计算机科学中最复杂的领域。而在那之前，\textbf{第三部分：抽象（Abstraction）}将率先解决你在此刻必然产生的重大疑虑：\enquote{程序员的工作难道就是日复一日地编写结构雷同的重复代码吗？}——我们将带你领略高阶函数与计算抽象的无穷魅力！
